\de{ĐỀ THI HỌC KỲ II NĂM HỌC 2022-2023}{THPT Nguyễn Thái Bình}

\Opensolutionfile{ans}[ans/ans-THPT-NguyenThaiBinh]


\begin{ex}%[0T0Y1-1]%[Dự án đề kiểm tra HKII NH22-23- Thy Nguyen Diem Vo]%[THPT Nguyễn Thái Bình]

	Ta quy ước SN có nghĩa là lần đầu tung mặt sấp, lần thứ hai là mặt ngửa, các kí hiệu SS, NS, NN được định nghĩa tương tự. Không gian mẫu khi tung một đồng xu hai lần là
	\choice
	{$\Omega=\{\mathrm{SS}, \mathrm{SN}\}$}
	{$\Omega=\{\mathrm{SS}, \mathrm{SN}, \mathrm{NS}\}$}
	{$\Omega=\{\mathrm{SS}\}$}
	{\True $\Omega=\{\mathrm{SS}, \mathrm{SN}, \mathrm{NS}, \mathrm{NN}\}$}
	\loigiai
	{
	Không gian mẫu khi tung một đồng xu hai lần là $\Omega=\{\mathrm{SS}, \mathrm{SN}, \mathrm{NS}, \mathrm{NN}\}$.
	}
\end{ex}
\begin{ex}%[0T8Y2-1]%[Dự án đề kiểm tra HKII NH22-23- Thy Nguyen Diem Vo]%[THPT Nguyễn Thái Bình]

	Số tập con gồm $k$ phần tử $(0 \leq k \leq n)$ lấy từ $n$ phần tử $(n \geq 1)$ của tập hợp $A$ là
	\choice
	{\True $\mathrm{C}_n^k$}
	{$n^k$}
	{$\mathrm{P}_n$}
	{$\mathrm{A}_n^k$}
	\loigiai
	{
	Số tập con gồm $k$ phần tử $(0 \leq k \leq n)$ lấy từ $n$ phần tử $(n \geq 1)$ của tập hợp $A$ là $\mathrm{C}_n^k$.	
	}
\end{ex}
\begin{ex}%[0T9Y2-1]%[Dự án đề kiểm tra HKII NH22-23- Thy Nguyen Diem Vo]%[THPT Nguyễn Thái Bình]

Cho các mệnh đề sau:
\begin{enumerate}
	\item Vectơ $\vec{n}$ khác $\overrightarrow{0}$ được gọi là vectơ pháp tuyến của đường thẳng $\Delta$ nếu giá của nó vuông góc với $\Delta$.
	\item Vectơ $\vec{u}$ khác $\overrightarrow{0}$ được gọi là vectơ chỉ phương của đường thẳng $\Delta$ nếu giá của nó song song với $\Delta$.
	\item Nếu $\vec{u}$ là vectơ chỉ phương của đường thẳng $\Delta$ thì $k \vec{u}$ $(k \neq 0)$ cũng là vectơ chỉ phương của $\Delta$.
	\item Nếu đường thẳng $\Delta$ có vectơ pháp tuyến $\vec{n}=(a ; b)$ thì $\Delta$ nhận $\vec{u}=(-b ;-a)$ là vectơ chỉ phương.
\end{enumerate}
Có bao nhiêu mệnh đề đúng?
	\choice
	{$1$}
	{\True $3$}
	{$2$}
	{$4$}
	\loigiai
	{
	Có 3 mệnh đề đúng là mệnh đề a, b, c.\\
	Mệnh đề d: Nếu đường thẳng $\Delta$ có vectơ pháp tuyến $\vec{n}=(a ; b)$ thì $\Delta$ nhận $\vec{u}=(b ;-a)$ là vectơ chỉ phương.
	}
\end{ex}
\begin{ex}%[0T9B4-2] %[Dự án đề kiểm tra HKII NH22-23- Thy Nguyen Diem Vo]%[THPT Nguyễn Thái Bình]

	Phương trình chính tắc của Elip có độ dài trục lớn bằng $12$ và trục nhỏ bằng $8$ là
	\choice
	{$\dfrac{x^2}{6}+\dfrac{y^2}{4}=1$}
	{$\dfrac{x^2}{36}-\dfrac{y^2}{16}=1$}
	{$\dfrac{x^2}{6}-\dfrac{y^2}{4}=1$}
	{\True $\dfrac{x^2}{36}+\dfrac{y^2}{16}=1$}
	\loigiai
	{
	Độ dài trục lớn bằng $12$, ta có $2a=12$ nên $a=6$.\\
	Độ dài trục nhỏ bằng $8$, ta có $2b=8$ nên $b=4$.\\
	Phương trình chính tắc của Elip là $\dfrac{x^2}{a^2}+\dfrac{y^2}{b^2}=1 \Leftrightarrow \dfrac{x^2}{36}+\dfrac{y^2}{16}=1$.
	}
\end{ex}
\begin{ex}%[0T7B2-1] %[Dự án đề kiểm tra HKII NH22-23- Thy Nguyen Diem Vo]%[THPT Nguyễn Thái Bình]

	Tam thức bậc hai $f(x)=2 x^2-3 x+1$ nhận giá trị dương khi
	\choice
	{$x \in(-\infty ; 1)$}
	{$x \in\left(\dfrac{1}{2} ;+\infty\right)$}
	{\True $x \in\left(-\infty ; \dfrac{1}{2}\right) \cup(1 ;+\infty)$}
	{$x \in\left(\dfrac{1}{2} ; 1\right)$}
	\loigiai
	{
	$2 x^2-3 x+1>0 \Leftrightarrow x \in\left(-\infty ; \dfrac{1}{2}\right) \cup(1 ;+\infty)$.	
	}
\end{ex}
\begin{ex}%[0T8Y3-1] %[Dự án đề kiểm tra HKII NH22-23- Thy Nguyen Diem Vo]%[THPT Nguyễn Thái Bình]
Khai triển nhị thức $(a+b)^4$, ta thu được kết quả đúng nào dưới đây?	
	\choice
	{$a^4-4 a^3 b+6 a^2 b^2-4 a b^3+b^4$}
	{$a^4+2 a^3 b+3 a^2 b^2+2 a b^3+b^4$}
	{$a^4+4 a^3 b+8 a^2 b^2+4 a b^3+b^4$}
	{\True $a^4+4 a^3 b+6 a^2 b^2+4 a b^3+b^4$}
	\loigiai
	{
	$(a+b)^4=a^4+4 a^3 b+6 a^2 b^2+4 a b^3+b^4$.
	}
\end{ex}
\begin{ex}%[0T7B1-1] %[Dự án đề kiểm tra HKII NH22-23- Thy Nguyen Diem Vo]%[THPT Nguyễn Thái Bình]
Đa thức nào là tam thức bậc hai? 
	\choice
	{\True $2 x^2-\sqrt{2} x+3$}
	{$x^3-2 x^2+1$}
	{$2x+1$}
	{$x^4+2 x^2$}
	\loigiai
	{
	$2 x^2-\sqrt{2} x+3$ là tam thức bậc hai.	
	}
\end{ex}
\begin{ex}%[0T9Y1-2]%[Dự án đề kiểm tra HKII NH22-23- Thy Nguyen Diem Vo]%[THPT Nguyễn Thái Bình]
Toạ độ của vectơ $\vec{a}=2 \vec{i}$ là
	\choice
	{$(0 ; 2)$}
	{\True $(2 ; 0)$}
	{$(2 ; 1)$}
	{$(1 ; 2)$}
	\loigiai
	{
		Toạ độ của vectơ $\vec{a}=2 \vec{i}$ là $2(1;0)=(2;0)$.
	}
\end{ex}
\begin{ex}%[0T8Y1-2] %[Dự án đề kiểm tra HKII NH22-23- Thy Nguyen Diem Vo]%[THPT Nguyễn Thái Bình]
	Một công việc được chia thành 2 công đoạn, công đoạn một có 2 cách thực hiện, công đoạn hai có 3 cách thực hiện. Khi đó, công việc có thể được thực hiện theo bao nhiêu cách?
	\choice
	{$4$}
	{\True $6$}
	{$3$}
	{$5$}
	\loigiai
	{
	Áp dụng quy tắc nhân, công việc có thể được thực hiện theo $2 \cdot 3=6$ cách.
	}
\end{ex}
\begin{ex}%[0T9Y1-1]%[Dự án đề kiểm tra HKII NH22-23- Thy Nguyen Diem Vo]%[THPT Nguyễn Thái Bình] 
	Cho $A(1;4)$, $B(-3; 2)$. Toạ độ trung điểm của đoạn thẳng $AB$ là
	\choice
	{$(-2 ; 3)$}
	{$(-1 ; 2)$}
	{$(-2 ; 1)$}
	{\True $(-1 ; 3)$}
	\loigiai
	{
 	Toạ độ trung điểm của đoạn thẳng $AB$ là $\left(\dfrac{1+(-3)}{2};\dfrac{4+2}{2}\right)=(-1;3)$.
	}
\end{ex}
\begin{ex}%[0T8Y2-1]%[Dự án đề kiểm tra HKII NH22-23- Thy Nguyen Diem Vo]%[THPT Nguyễn Thái Bình]
	Số các hoán vị của $n$ phần tử là
	\choice
	{$n$}
	{$n\cdot n$}
	{$n^n$}
	{\True $\mathrm{P}_n$}
	\loigiai
	{
		Số các hoán vị của $n$ phần tử là $\mathrm{P}_n=n!$.
	}
\end{ex}
\begin{ex}%[0T7Y2-1]%[Dự án đề kiểm tra HKII NH22-23- Thy Nguyen Diem Vo]%[THPT Nguyễn Thái Bình]
	Tập nghiệm của bất phương trình bậc hai $x^2-4x+3\ge 3$ là
	\choice
	{\True $S=(-\infty;0]\cup [4;+\infty)$}
	{$S=(-\infty;0)\cup (4;+\infty)$}
	{$S=(-\infty;1]\cup [3;+\infty)$}
	{$S=(-\infty;1)\cup (3;+\infty)$}
	\loigiai
	{
		Ta có $x^2-4x+3\ge 3\Leftrightarrow x^2-4x\ge 0$.\\
		Đặt $f(x)=x^2-4x$. Ta có bảng xét dấu
		\begin{center}
			\begin{tikzpicture}
				\tkzTabInit[lgt=2,espcl = 3]
				{$x$ /0.7, $f(x)$ /0.7}
				{$-\infty$,$0$,$4$,$+\infty$}
				\tkzTabLine{ ,+,$0$,-,$0$, +, }
			\end{tikzpicture}
		\end{center}
		Từ bảng xét dấu, ta có tập nghiệm $S=(-\infty;0]\cup [4;+\infty)$.
	}
\end{ex}
\begin{ex}%[0T8Y1-1]%[Dự án đề kiểm tra HKII NH22-23- Thy Nguyen Diem Vo]%[THPT Nguyễn Thái Bình]
	Một công việc được thực hiện theo phương án $A$ hoặc phương án $B$, phương án $A$ có $5$ cách thực hiện, phương án $B$ có $7$ cách thực hiện không trùng với bất kỳ cách nào của phương án $A$. Khi đó, công việc có thể được thực hiện theo bao nhiêu cách?
	\choice
	{$36$}
	{\True $12$}
	{$35$}
	{$11$}
	\loigiai
	{
		Theo quy tắc cộng, công việc có thể được thực hiện theo $7+5=12$ cách.
	}
\end{ex}
\begin{ex}%[0T0Y1-2]%[Dự án đề kiểm tra HKII NH22-23- Thy Nguyen Diem Vo]%[THPT Nguyễn Thái Bình]
	Cho $A$ là biến cố và $\overline{A}$ là biến cố đối của biến cố A, trong các mệnh đề sau, mệnh đề nào đúng?
	\choice
	{\True $\mathrm{P}(A)+\mathrm{P}(\overline{A})=1$}
	{$\mathrm{P}(A)\cdot\mathrm{P}(\overline{A})=1$}
	{$\mathrm{P}(A)-\mathrm{P}(\overline{A})=1$}
	{$\mathrm{P}(\overline{A})=\mathrm{P}(A)$}
	\loigiai
	{
		Mệnh đề $\mathrm{P}(A)+\mathrm{P}(\overline{A})=1$ đúng.
	}
\end{ex}
\begin{ex}%[0T9Y3-2]%[Dự án đề kiểm tra HKII NH22-23- Thy Nguyen Diem Vo]%[THPT Nguyễn Thái Bình]
	Đâu là phương trình đường tròn $(C)$ có tâm $I(a;b)$ và bán kính bằng $R$?
	\choice
	{$(x+a)^2+(y+b)^2=R^2$}
	{$(x-a)^2-(y-b)^2=R^2$}
	{\True $(x-a)^2+(y-b)^2=R^2$}
	{$(x-a)^2+(y-b)^2=R$}
	\loigiai
	{
		Phương trình đường tròn $(C)$ có tâm $I(a;b)$ và bán kính bằng $R$ là $(x-a)^2+(y-b)^2=R^2$.
	}
\end{ex}
\begin{ex}%[0T0Y2-1]%[Dự án đề kiểm tra HKII NH22-23- Thy Nguyen Diem Vo]%[THPT Nguyễn Thái Bình]
	$A$ là biến cố của phép thử có không gian mẫu $\Omega $, trong các mệnh đề sau, mệnh đề nào \textbf{sai}?
	\choice
	{$\mathrm{P}(\Omega)=1$}
	{\True $0<\mathrm{P}(A)<1$}
	{$\mathrm{P}(\varnothing)=0$}
	{$\mathrm{P}(A)=\dfrac{n(A)}{n(\Omega)}$}
	\loigiai
	{
		Các mệnh đề $\mathrm{P}(\Omega)=1$, $\mathrm{P}(\varnothing)=0$, $\mathrm{P}(A)=\dfrac{n(A)}{n(\Omega)}$ đúng. Mệnh đề $0<\mathrm{P}(A)<1$ sai vì với mọi biến cố $A$, ta có $0\le\mathrm{P}(A)\le 1$.
	}
\end{ex}
\begin{ex}%[0T8Y2-1]%[Dự án đề kiểm tra HKII NH22-23- Thy Nguyen Diem Vo]%[THPT Nguyễn Thái Bình]
	Một tổ học sinh gồm có $6$ nam và $4$ nữ. Có bao nhiêu cách chọn $5$ học sinh của tổ tham gia đội xung kích?
	\choice
	{\True $\mathrm{C}_{10}^5$}
	{$\mathrm{C}_6^3\cdot\mathrm{C}_7^2$}
	{$\mathrm{A}_{10}^5$}
	{$\mathrm{C}_6^4\cdot\mathrm{C}_8^1$}
	\loigiai
	{
		Số cách chọn $5$ học sinh của tổ tham gia đội xung kích là $\mathrm{C}_{10}^5$.
	}
\end{ex}
\begin{ex}%[0T9Y1-2]%[Dự án đề kiểm tra HKII NH22-23- Thy Nguyen Diem Vo]%[THPT Nguyễn Thái Bình]
	Cho hai véc-tơ $\overrightarrow{a}=(1;-4)$ và $\overrightarrow{b}=(-3;5)$. Toạ độ của véc-tơ $2\overrightarrow{a}+\overrightarrow{b}$ là
	\choice
	{\True $(-1;-3)$}
	{$(2;1)$}
	{$(1;-3)$}
	{$(-2;-1)$}
	\loigiai
	{
		Ta có $2\overrightarrow{a}=(2;-8)$, suy ra $2\overrightarrow{a}+\overrightarrow{b}=(-1;-3)$.
	}
\end{ex}
\begin{ex}%[0T7B3-2]%[Dự án đề kiểm tra HKII NH22-23- Thy Nguyen Diem Vo]%[THPT Nguyễn Thái Bình]
	Tập nghiệm của phương trình $\sqrt{x^2-3x+2}=\sqrt{3x-7}$ là
	\choice
	{$S=\{4\}$}
	{\True $S=\{3\}$}
	{$S=\{4;2\}$}
	{$S=\{1;3\}$}
	\loigiai
	{
		Bình phương hai vế của phương trình, ta được
		\[x^2-3x+2=3x-7\Leftrightarrow x^2-6x+9=0\Leftrightarrow x=3. \]
		Thử lại, ta thấy $x=3$ thỏa mãn. Vậy tập nghiệm của phương trình là $S=\{3\}$.
	}
\end{ex}
\begin{ex}%[0T7B1-1]%[Dự án đề kiểm tra HKII NH22-23- Thy Nguyen Diem Vo]%[THPT Nguyễn Thái Bình]
	Cho bảng xét dấu của một tam thức bậc hai bên dưới sau đây
	\begin{center}
		\begin{tikzpicture}[line join = round, line cap = round,>=stealth,scale=1]
			\tkzTabInit[lgt=1.2,espcl=3,deltacl=0.5]
			{$x$ /0.7,  $f(x)$ /0.7}
			{$-\infty$,$-3$,$3$,$+\infty$}
			\tkzTabLine{ ,-,$0$,+,$0$,-, }
		\end{tikzpicture}
	\end{center}
	Mệnh đề nào sau đây là đúng?
	\choice
	{$f(x)=x^2+x-6$}
	{\True $f(x)\ge 0\Leftrightarrow x\in [-3;2]$}
	{$f(x)\ge 0\Leftrightarrow x\in (-\infty;2)$}
	{$f(x)<0\Leftrightarrow x\in (-3;+\infty)$}
	\loigiai
	{
		Từ bảng xét dấu, ta thấy $f(x)\ge 0\Leftrightarrow x\in [-3;2]$.
	}
\end{ex}

\begin{ex}%[0T0Y2-2]%[Dự án đề kiểm tra HKII NH22-23- Nguyen Huynh]%[THPT Nguyễn Thái Bình]
	Gieo hai con xúc xắc cân đối và đồng chất. Xác suất để tổng số chấm trên hai mặt xuất hiện bằng $7$ là 
	\choice 
	{$\dfrac{1}{3}$}
	{$\dfrac{2}{3}$}
	{\True $\dfrac{1}{6}$}
	{$\dfrac{1}{9}$}
	\loigiai{
		Không gian mẫu $n(\Omega)=6\cdot6=36$\\
		Biến cố ${A} \colon$ \lq\lq Tổng số chấm trên hai mặt xuất hiện bằng $7$ \rq\rq.\\
		Các trường hợp xảy ra của biến cố $A$ là $\{1,6\}$, $\{6,1\}$, $\{2,5\}$, $\{5,2\}$, $\{3,4\}$, $\{4,3\}$.\\
		Xác xuất của biến cố $A$  $\mathrm {P(A)}=\dfrac{n({A})}{n(\Omega)}=\dfrac{6}{36}$.
	}
\end{ex}

\begin{ex}%[0T9Y3-1]%[Dự án đề kiểm tra HKII NH22-23- Nguyen Huynh]%[THPT Nguyễn Thái Bình]
	Phương trình đường tròn $(C) \colon x^2+y^2-2x+4y-6=0$ có tâm là
	\choice 
	{$(-1;2)$}
	{\True $(1;-2)$}
	{$(-1;-2)$}
	{$(1;2)$}
	\loigiai{
		Đường tròn $(C) \colon x^2+y^2-2ax-2by+c=0$ có tâm $I(a;b)$.\\
		Do đó đường tròn $(C) \colon x^2+y^2-2x+4y-6=0$ có tâm là $I(1;-2)$.
	}
\end{ex}

\begin{ex}%[0T8Y3-2]%[Dự án đề kiểm tra HKII NH22-23- Nguyen Huynh]%[THPT Nguyễn Thái Bình]
	Hệ số của $x^3$ trong khai triển $(2x-3)^4$ là
	\choice 
	{\True $-96$}
	{$96$}
	{$-216$}
	{$216$}
	\loigiai{Ta có
		\allowdisplaybreaks
		\begin{eqnarray*}
			(2x-3)^4 &=&(2x)^4-4\cdot(2x)^3\cdot3+ 6\cdot(2x)^2\cdot3^2- 4\cdot(2x)\cdot3^3 + 3^4\\
			&=& 16x^4 -96x^3 +216x^2- 216x+ 81
		\end{eqnarray*}	
		Vậy hệ số của $x^3$ trong khai triển $(2x-3)^4$ là $-96$.
		
	}
\end{ex}

\begin{ex}%[0T8B2-1]%[Dự án đề kiểm tra HKII NH22-23- Nguyen Huynh]%[THPT Nguyễn Thái Bình]
	Có bao nhiêu số tự nhiên có $5$ chữ số (các chữ số không nhất thiết khác nhau) lấy từ tập $X=\{1,2,3\}$?
	\choice 
	{$\mathrm {A}_5^3$}
	{\True $3^5$}
	{$5^3$}
	{$\mathrm {C}_5^3$}
	\loigiai{
		Gọi $\overline{abcde}$ là số có ba chữ số cần tìm.
		\begin{itemize}
			\item $a$ có $3$ cách chọn.
			\item $b$ có $3$ cách chọn.
			\item $c$ có $3$ cách chọn.
			\item $d$ có $3$ cách chọn.
			\item $e$ có $3$ cách chọn.
		\end{itemize}
		Vậy có tổng cộng $3^5$ số tự nhiên cần tìm.}
\end{ex}

\begin{ex}%[0T8Y1-2]%[Dự án đề kiểm tra HKII NH22-23- Nguyen Huynh]%[THPT Nguyễn Thái Bình]
	Một người có $6$ cái quần khác nhau, $5$ cái áo khác nhau, $4$ chiếc cà vạt khác nhau. Để chọn một cái quần hoặc một cái áo hoặc một cái cà vạt thì số cách chọn khác nhau là
	\choice 
	{$74$}
	{\True $120$}
	{$455$}
	{$15$}
	\loigiai{
		\begin{itemize}
			\item Có $6$ cách chọn cái quần khác nhau.
			\item Có $5$ cách chọn cái áo khác nhau.
			\item Có $4$ cách chọn chiếc cà vạt khác nhau.
		\end{itemize}	
		Vậy có $6 \cdot 5 \cdot 4 = 120$ cách chọn một cái quần hoặc một cái áo hoặc một cái cà vạt khác nhau.}
\end{ex}

\begin{ex}%[0T9B2-2]%[Dự án đề kiểm tra HKII NH22-23- Nguyen Huynh]%[THPT Nguyễn Thái Bình]
	Phương trình tham số của đường thẳng đi qua 2 điểm $M(1;-2)$, $N(-3;2)$ là
	\choice 
	{\True $\heva{&x=-3+t\\&y=2+t}$}
	{$\heva{&x=1-t \\&y=-2+t}$}
	{$\heva{&x=-1+t\\&y=1-2t}$}
	{$\heva{&x=1+4t\\&y=-2+4t}$}
	\loigiai{
		Phương trình tham số của đường thẳng $MN$ đi qua điểm $N(-3;2)$ và nhận véc-tơ $\overrightarrow{MN}=(-4;4)$ hay $\vec{u}=(1;1)$ làm véc-tơ chỉ phương là\\
		$$(MN) \colon \heva{&x=-3+t \\ &y=2+t.}$$
	}
\end{ex}
\begin{ex}%[0T0B1-1]%[Dự án đề kiểm tra HKII NH22-23- Nguyen Huynh]%[THPT Nguyễn Thái Bình]
	Một hộp chứa 6 bi xanh và 8 bi đỏ, lấy ngẫu nhiên cùng một lúc từ trong hộp ra 4 bi. Số phần tử của không gian mẫu là
	\choice 
	{ $\mathrm C_8^4$} 
	{ $\mathrm A_{14}^4$} 
	{\True $\mathrm C_{14}^4$} 
	{ $\mathrm C_6^4$}
	\loigiai{
		Không gian mẫu là tập hợp các cách lấy ra 4 viên bi trong $14$ viên bi.\\
		Do đó, số phần tử của không gian mẫu là $\mathrm C_{14}^4$.
	}
\end{ex} 
\begin{ex} %[0T7Y1-1]%[Dự án đề kiểm tra HKII NH22-23- Nguyen Huynh]%[THPT Nguyễn Thái Bình] 
	Với giá trị nào của $m$ thì đa thức $(m-1) x^2+2 m x+3$ \textbf{không} là tam thức bậc hai?
	\choice 
	{ $m=2$} 
	{\True $m=1$} 
	{ $m=-1$} 
	{ $m=0$}
	\loigiai{
		Đa thức trên không là tam thức bậc hai khi $a=0 \Leftrightarrow m-1=0 \Leftrightarrow m=1$.
	}
\end{ex} 
\begin{ex} %[0T8K2-1]%[Dự án đề kiểm tra HKII NH22-23- Nguyen Huynh]%[THPT Nguyễn Thái Bình] 
	Có bao nhiêu số tự nhiên chẵn có $5$ chữ số khác nhau?
	\choice 
	{\True $13776$} 
	{ $13440$} 
	{ $13600$} 
	{ $10752$}
	\loigiai{
		Kí hiệu số cần lập có dạng $\overline{abcde}$.
		\begin{itemize}
			\item Trường hợp 1: $e=0$.\\
			Chọn $\overline{abcd}$ có $\mathrm A^4_9$ cách.
			\item Trường hợp 2: $e \in \{2;\;4;\;6;\;8\}$ có 4 cách chọn.
			\begin{itemize}
				\item Chọn $a\;(a\ne 0, a\ne e)$ có 8 cách chọn.
				\item Chọn $\overline{bcd}$ có $\mathrm A^3_8$ cách chọn.
			\end{itemize}
		\end{itemize}
		Vậy có $\mathrm A^4_9+4\cdot8\cdot \mathrm A^3_8=13776$ số chẵn có 5 chữ số khác nhau.
	}
\end{ex} 
\begin{ex} %[0T7K2-2]%[Dự án đề kiểm tra HKII NH22-23- Nguyen Huynh]%[THPT Nguyễn Thái Bình] 
	Một vật được ném theo phương thẳng đứng xuống dưới từ độ cao $120 \mathrm m$ với vận tốc ban đầu $v_0=20 \mathrm{m / s}$. Hỏi sau ít nhất bao nhiêu giây, vật đó cách mặt đất không quá $50\mathrm m$? Giả thiết rằng sức cản của không khí là không đáng kể, lấy $g=10 \mathrm {m/s^2}$?
	\choice 
	{ $2,25$} 
	{ $2,5$} 
	{\True $2,24$} 
	{ $3,5$}
	\loigiai{
		Ta có quãng đường của vật đi được sau khoảng thời gian $t$ giây là $s=v_0t+\dfrac{a t^2}{2}=20t+5t^2$ với $t>0$.\\
		Khi đó, độ cao của vật là $h=120-s=120-20t-5t^2$.\\
		Khi đó, $120-20t-5t^2\ge 50 \Leftrightarrow t\ge -2+3\sqrt2 \approx 2,24 $.
		
	}
\end{ex} 
\begin{ex} %[0T9B4-0]%[Dự án đề kiểm tra HKII NH22-23- Nguyen Huynh]%[THPT Nguyễn Thái Bình] 
	Mặt trăng di chuyển quanh Trái đất theo quỹ đạo hình elip với Trái đất là một tiêu điểm. Các trục lớn và trục bé của quỹ đạo có độ dài lần lượt là $768800 \mathrm{~km}$ và $767640 \mathrm{~km}$. Khoảng cách lớn nhất từ tâm Trái Đất đến tâm Mặt trăng xấp xỉ bao nhiêu km?
	\choice 
	{ $420608$} 
	{\True $405508$} 
	{ $412615$} 
	{ $430006$}
	\loigiai{
		Ta có $2a=768800 \Rightarrow a=384400,\;2b=767640 \Rightarrow b=383820$.\\
		Khi đó $c=\sqrt{a^2-b^2}\approx 21108$.\\
		Vậy khoảng cách lớn nhất từ Trái đất đến Mặt trăng là $a+c\approx 405508 \mathrm{(km)}$.
} \end{ex} 
\begin{ex} %[0T7B3-2]%[Dự án đề kiểm tra HKII NH22-23- Nguyen Huynh]%[THPT Nguyễn Thái Bình]
	Tổng các nghiệm của phương trình $\sqrt{2 x^2-6 x+3}=x+1$ bằng
	\choice 
	{ 6} 
	{ 4} 
	{\True 8} 
	{ 7 }
	\loigiai{
		$$\begin{array}{ll}
			&\sqrt{2 x^2-6 x+3}=x+1 \\
			\Leftrightarrow & \heva{&x+1\ge 0\\&2x^2-6x+3=(x+1)^2}\\
			\Leftrightarrow & \heva{&x\ge -1 \\& x^2-8x+2=0}\\
			\Leftrightarrow & \heva{&x\ge -1\\& \hoac{&x=4+\sqrt{14}\\& x=4-\sqrt{14}}}\\
			\Leftrightarrow & \hoac{&x=4+\sqrt{14}\\& x=4-\sqrt{14}}\\
		\end{array}$$
		Vậy tổng hai nghiệm của phương trình là $8$.
	}
\end{ex}

\begin{ex}%[0T8B3-2]%[Dự án đề kiểm tra HKII NH22-23- Nguyen Huynh]%[THPT Nguyễn Thái Bình]
	Khai triển biểu thức $(a+b x)^4$ theo thứ tự bậc của $x$ tăng dần, ta nhận được biểu thức gồm hai số hạng đầu tiên theo thứ tự là $ 1 $ và $-8 x$. Tích của $a$ và $b$ bằng 
	\choice
	{$4$}
	{\True $-2$}
	{$-4$}
	{$2$}
	\loigiai{
		Ta có $a^{4}=1$ và $4\cdot a^{3}b=-8$, suy ra $ab=-2$.
	}
\end{ex}

\begin{ex}%[0T9B2-2]%[Dự án đề kiểm tra HKII NH22-23- Nguyen Huynh]%[THPT Nguyễn Thái Bình] 
	Cho tam giác $ABC$ có $A(2 ; 8)$, $B(-1 ; 2)$, $C(5 ; 4)$. Phương trình đường trung tuyến kẻ từ $A$ là
	\choice
	{$y=2$}
	{\True $x=2$}
	{$y=-2$}
	{$x=-2$}
	\loigiai{
		Gọi $M$ là trung điểm $BC$, suy ra $M(2;3) \Rightarrow \overrightarrow{AM}=(0;-5)$.
		\\ Chọn $\vec{n}_{AB} = (1;0)$, suy ra $$d\colon 1(x-2)+0(y-3)=0 \Leftrightarrow x=2$$
	}
\end{ex}

\begin{ex}%[0T9B3-3]%[Dự án đề kiểm tra HKII NH22-23- Nguyen Huynh]%[THPT Nguyễn Thái Bình]
	Phương trình tiếp tuyến của đường tròn $(C)\colon (x-3)^2+(y+2)^2=25$ tại điểm có hoành độ bằng $ 6 $ và tung độ dương là
	\choice
	{\True $3 x+4 y-26=0$}
	{$3 x+4 y-16=0$}
	{$3 x+4 y-27=0$}
	{$3 x-4 y+26=0$}
	\loigiai{
		Thay $x=6$ vào phương trình của $(C)\colon (6-3)^{2} + (y+2)^{2} = 15 \Rightarrow y = 2 $ (do $y_{M}>0$).\\ Ta có $M(6;2)$ và tâm $I(3;-2)$, suy ra $\overrightarrow{IM} = (3;4)$.\\ Vậy phương trình tiếp tuyến  $d\colon 3(x-6)+4(y-2)=0 \Leftrightarrow 3x+4y-26=0$.
	}
\end{ex}

\begin{ex}%[0T8B2-1]%[Dự án đề kiểm tra HKII NH22-23- Nguyen Huynh]%[THPT Nguyễn Thái Bình]
	Một nhóm gồm $ 12 $ bạn đăng ký tham gia các môn thể thao, biết rằng trong danh sách cần $ 6 $ bạn chơi đá banh, $ 4 $ bạn chơi bóng rổ và $ 2 $ bạn chơi bóng chuyền. Số cách lập ra một danh sách trên là:
	\choice
	{$14580$}
	{$12534$}
	{\True $13860$}
	{$16720$}
	\loigiai{
		Số cách chọn 6 bạn (từ 12 bạn) chơi đá banh là: $\mathrm C_{12}^{6}$ (cách)\\
		Số cách chọn 4 (từ 6 bạn) chơi bóng rổ là: $\mathrm C_{6}^{4}$ (cách)\\
		Số cách chọn 2 bạn chơi bóng chuyền là: $1$ (cách)\\
		Số cách lập ra một danh sách thỏa mãn là: $\mathrm C_{12}^{6} \cdot \mathrm C_{6}^{4} \cdot 1 = 13860$.
	}
\end{ex}

\begin{ex}%[0T7B1-2]%[Dự án đề kiểm tra HKII NH22-23- Nguyen Huynh]%[THPT Nguyễn Thái Bình]
	Gọi $S$ là tập hợp các giá trị của $m$ thỏa mãn bất phương trình $x^2-2 m x+4 \geq 0, \forall x \in \mathbb{R}$. Khi đó
	\choice
	{$S=(-2 ; 2)$}
	{$S=(2 ;+\infty)$}
	{$S=(-\infty ;-2] \cup[2 ;+\infty)$}
	{\True $S=[-2 ; 2]$}
	\loigiai{
		Yêu cầu bài toán tương đương $\Delta^{\prime} = m^{2} - 4 \leq 0 \Leftrightarrow -2 \leq m \leq 2$.
	}
\end{ex}

\begin{ex}%[0T9K3-6]%[Dự án đề kiểm tra HKII NH22-23- Nguyen Huynh]%[THPT Nguyễn Thái Bình]
	Một khu vườn hình chữ nhật $A B C D$ có diện tích bằng $20 \text{m}^2$, biết rằng các đường thẳng $A B$, $B C, D A$ lần lượt có phương trình là : $y=5, x=7, x=2$. Người ta muốn lắp đặt một vòi phun nước xoay tròn để tưới nước cho khu vườn thì bán kính phun tối thiểu của vòi để tưới hết khu vườn nằm trong khoảng :
	\choice
	{$(3,5 ; 4)$}
	{\True $(3 ; 3,5)$}
	{$(4 ; 4,5)$}
	{$(2,5 ; 3)$}
	\loigiai{
		Từ giả thiết suy ra $AB = 5, \, AD = 4$. Vậy $R_{min} = \dfrac{1}{2} \sqrt{4^{2} + 5^{2}} \approx 3,2$ (m).
	}
\end{ex}

\begin{ex}%[0T8K2-2]%[Dự án đề kiểm tra HKII NH22-23- Nguyen Huynh]%[THPT Nguyễn Thái Bình]
	Trong hộp có 5 bi xanh, 7 bi đỏ và 10 bi vàng có kích thước và khối lượng như nhau. Lấy ngẫu nhiên từ trong hộp ra 8 viên bi. Xác suất để trong 8 bi lấy ra có không quá 4 bi xanh là:
	\choice
	{$\dfrac{208}{209}$}
	{$\dfrac{514}{627}$}
	{$\dfrac{1544}{1881}$}
	{\True $\dfrac{1877}{1881}$}
	\loigiai{
		Ta có số phần tử của không gian mẫu $n(\Omega) = \mathrm C_{22}^{8}$.\\
		Gọi $A$ là biến cố mà "Trong 8 bi lấy ra có không quá 4 bi xanh".\\
		Suy ra $\overline{A}$: \lq\lq Trong 8 bi lấy ra có 5 bi xanh\rq\rq.Suy ra $n(\overline{A}) = \mathrm C_{5}^{5} \cdot \mathrm C_{17}^{3}$.\\
		Vậy $P(A) = 1 - \dfrac{\mathrm C_{5}^{5}\cdot \mathrm C_{17}^{3}}{\mathrm C_{22}^{8}} = \dfrac{1877}{1881}$.
	}
\end{ex}

\begin{ex}%[0T8K2-1]%[Dự án đề kiểm tra HKII NH22-23- Nguyen Huynh]%[THPT Nguyễn Thái Bình]
	Cần sắp xếp thứ tự $ 10 $ tiết mục văn nghệ cho buổi biển diễn văn nghệ của trường. Ban tổ chức dự kiến xếp $ 4 $ tiết mục ca nhạc ở vị trí thứ $ 1 $, thứ $ 3 $, thứ $ 4 $ và thứ $ 9 $; xếp $ 4 $ tiết mục múa ở vị trí thứ $ 2 $ , thứ $ 5 $ , thứ $ 7 $ và thứ $10$ ; xếp $2$ tiết mục kịch ở vị trí thứ $ 6 $ và thứ $ 8 $. Có bao nhiêu cách xếp khác nhau?
	\choice
	{$1856$}
	{$1800$}
	{$1480$}
	{\True $1152$}
	\loigiai{
		Xếp 4 tiết mục ca nhạc có $4!$ cách.\\
		Xếp 4 tiết mục múa có $4!$ cách.\\
		Xếp 2 tiết mục kịch có $2!$ cách.\\
		Số cách xếp khác nhau là $4! \cdot 4! \cdot 2! = 1152$.
	}
\end{ex}

\Closesolutionfile{ans}
%\begin{center}
%	\textbf{ĐÁP ÁN}
%	\inputansbox{10}{ans/ans-THPT-NguyenThaiBinh}	
%\end{center}
